\documentclass[10pt, a4paper]{article}
\usepackage[polish]{babel}
\usepackage[utf8]{inputenc}
\usepackage{microtype}
\usepackage[table]{xcolor}
\usepackage{fontspec}
\usepackage{setspace}
\usepackage{fancyhdr}
\usepackage{tabularx}
\usepackage{geometry}
\usepackage{colortbl}

\renewcommand{\headrulewidth}{0pt}
\renewcommand\tabularxcolumn[1]{m{#1}}

\setmainfont{Ubuntu-Regular.ttf}[
  BoldFont = Ubuntu-Bold.ttf,
  ItalicFont = Ubuntu-Italic.ttf,
  BoldItalicFont = Ubuntu-BoldItalic.ttf,
]

\newgeometry{
    top=25mm,
    bottom=25mm,
    outer=15mm,
    inner=15mm
}

\setstretch{1.1}

\pagestyle{fancy}
\fancyhf{}
\fancyfoot[C]{\footnotesize str. \thepage}

\begin{document}
\sloppy

\begin{center}
{\Large\bfseries Formularz Wyników Drużyny}\\[12pt]
{\normalsize
\textbf{Drużyna:} <%teamName%> \\
\textbf{Miasto:} <%cityName%> \\
\textbf{Problem:} <%problem%> \hspace{1em} \textbf{Grupa wiekowa:} <%age%> \\
<%#performanceAt%>\textbf{Godzina występu:} <%performanceAt%><%/performanceAt%><%#performanceTime%> \hspace{1em} \textbf{Czas trwania:} <%performanceTime%> min<%/performanceTime%>
}\\[16pt]
\end{center}

\noindent\textbf{Długoterminowe (DT)}\\[4pt]
\begin{tabularx}{\linewidth}{|X|r|}
\hline
\rowcolor{gray!15}
\textbf{Wpis} & \textbf{Średnia} \\
\hline
<%#dtEntries%>
<%#isSectionHeader%>
\textbf{<%name%>} & --- \\
<%/isSectionHeader%>
<%^isSectionHeader%><%#noElement%>
<%name%> & \textit{brak elementu} \\
<%/noElement%><%^noElement%>
<%name%> & <%score%> \\
<%/noElement%><%/isSectionHeader%>
<%/dtEntries%>
\hline
\rowcolor{blue!10}
\textbf{Suma DT} & \textbf{<%dtSum%>} \\
\hline
\end{tabularx}

\vspace{12pt}
\noindent\textbf{Styl}\\[4pt]
\begin{tabularx}{\linewidth}{|X|r|}
\hline
\rowcolor{gray!15}
\textbf{Wpis} & \textbf{Średnia} \\
\hline
<%#styleEntries%>
<%name%> & <%score%> \\
<%/styleEntries%>
\hline
\rowcolor{blue!10}
\textbf{Suma Styl} & \textbf{<%styleSum%>} \\
\hline
\end{tabularx}

\vspace{12pt}
\noindent\textbf{Kary}\\[4pt]
\begin{tabularx}{\linewidth}{|X|r|}
\hline
\rowcolor{gray!15}
\textbf{Wpis} & \textbf{Wynik} \\
\hline
<%#penaltyEntries%>
<%name%> & <%score%> \\
<%/penaltyEntries%>
\hline
\rowcolor{blue!10}
\textbf{Suma Kar} & \textbf{<%penaltySum%>} \\
\hline
\end{tabularx}

\vspace{12pt}
\begin{tabularx}{\linewidth}{|X|r|}
\hline
\rowcolor{green!15}
\textbf{SUMA CAŁKOWITA} & \textbf{<%totalSum%>} \\
\hline
\end{tabularx}

\end{document}
